\chapter{The Problem of Emergent Fermions}

The framework has matured to a level where we no longer need to question whether it can produce physics - it clearly can. 

However, there were no stable collisions, no bouncing, and no persistent separation.

We had successfully emergent bosons. We did not have emergent fermions (half-integer spin particles that
obey the Pauli exclusion principle).

\subsection*{The Core Difficulty}

In standard quantum mechanics, the exchange of two identical particles produces a characteristic phase:
\begin{itemize}
    \item Bosons: \(\psi(x_1, x_2) = +\psi(x_2, x_1)\)
    \item Fermions: \(\psi(x_1, x_2) = -\psi(x_2, x_1)\)
\end{itemize}

The minus sign forces the wavefunction to vanish whenever \(x_1 = x_2\), creating a strict exclusion effect.

Our current spectral complexity measure, which counts the informational cost of frequencies, amplitudes, and phases, naturally favors symmetric (bosonic) configurations.
Stacking multiple identical excitations on the same modes is cheaper — one description effectively covers many particles.
This leads to merging behavior, as symmetric states minimize total description length.

In contrast, producing a sharp antisymmetric cancellation that drives the probability exactly to zero at overlap points appears expensive in the current measure.
It seems to require either many high-frequency components (heavily suppressed) or an additional rule that we have not yet derived from pure informational cost.

Without some form of emergent exclusion, localized coherent structures cannot maintain stable individual identities over time.
Atoms cannot form. Chemistry becomes impossible. Complex matter — and therefore observers — cannot exist in any stable form.


\section{Where to Go From Here}

Hard-coding antisymmetry, or adding it as an extra postulate, would violate the central goal: everything must emerge from the self-interpretive informational dynamics and the induced measure.
The only path consistent with the spirit of this framework is to find a purely informational mechanism by which antisymmetric configurations become dramatically lower in total joint description length
than their symmetric counterparts.

\subsection*{Intuition}

When two observer worldlines overlap in space while remaining separate in their full 4D extent, the symmetric (bosonic) description inevitably forces a merging of identities across time.
The resulting 4D structure becomes highly disordered in the overlap region, requiring either:

\begin{itemize}
  \item a massive injection of high-frequency components to track the chaos, or
  \item a loss of distinct observer identities altogether.
\end{itemize}

Both options produce a catastrophic drop in the induced observer measure.

So some antisymmetric sign-flip cost is needed. Although this adds some local descriptive overhead, it is still cheaper and gets automatically selected.
The minus sign is then not imposed by hand but emerges as the smaller cost.



\subsection*{Refining Spectral Complexity}

Kolmogorov complexity was too discrete. The current, first-order version of Spectral Complexity — which adds the bit-cost of individual frequencies, amplitudes, and phases — naturally favors bosonic behavior is apparently
too trivial and naive. We need a deeper, more sophisticated measure.

The true informational cost of a wavefunction must include not only the direct parameters of its modes, but also the cost of the \emph{encoding scheme} chosen to represent those modes.

The \nameref{pr:siip} must be applied recursively to the spectral compression itself.

If the theory is correct, the sudden drop in complexity should pick up the "Fermionic encoding", which will cause a measure-rebound.


\section{Proof of Concept}

[ UNDER CONSTRUCTION ]

